\documentclass[11pt,letterpaper]{article}

\usepackage[T1]{fontenc}
\usepackage{amsmath,amssymb}
\usepackage{tikz}
\usepackage{pgfplots}
\pgfplotsset{compat=1.18}
\usepackage{algorithm}
\usepackage{algpseudocode}
\usepackage{booktabs}
\usepackage{enumitem}
\usepackage[hidelinks]{hyperref}

\newcommand{\N}{\mathbb{N}}
\newcommand{\Z}{\mathbb{Z}}
\newcommand{\Q}{\mathbb{Q}}
\newcommand{\R}{\mathbb{R}}
\newcommand{\C}{\mathbb{C}}
\newcommand{\Map}[3]{#1 \colon #2 \to #3}
\newcommand{\Mod}[1]{#1\text{-}\mathrm{Mod}}
\newcommand{\GroupHom}{\mathrm{Hom}_{\mathbf{Grp}}}
\newcommand{\RingHom}{\mathrm{Hom}_{\mathbf{Ring}}}

\newtheorem{theorem}{Theorem}
\newtheorem{lemma}[theorem]{Lemma}
\newtheorem{definition}[theorem]{Definition}
\newtheorem{remark}[theorem]{Remark}
\newenvironment{proof}{\noindent\textit{Proof.}}{\hfill\(\square\)}

\title{Generalized Collatz Library}
\author{ascendnoosphere}
\date{\today}

\begin{document}
\maketitle

\begin{abstract}
A multiplier-supported divisor of the offset gives an absorbing integer
sublattice. A lexicographic descent bounds the affine instructions before
entry, and a single entry witness determines the entire normalized raw tail.
Every periodic state already belongs to the sublattice; every cycle phase
and primitive raw period recover by scaling. Repeated gcd removal constructs
the maximal supported factor without prime factorization and leaves an
offset coprime to the unchanged multiplier. The resulting search cutoff is
an exact ceiling, not a rounded-down approximation. These results concern
normalization and recovery, not convergence of the remaining core map.
\end{abstract}

\section{The supplied factor and its invariant lattice}

Let \(Q,C,D\in\Z\) with \(|D|\ge2\), and define the raw map
\[
F(n)=\begin{cases}
n/D,&D\mid n,\\
Qn+C,&D\nmid n.
\end{cases}
\]
Division has precedence, including at zero, and retains the sign of \(D\).
The absorption identity is established below. The GCD construction and the
search-cutoff conversion are also developed here; coprime scaling at the
same raw clock appears in \cite{AscendScaling}.

\begin{definition}
A supplied factor is an integer \(S\) satisfying
\[
S\mid C,\qquad \gcd(|S|,|D|)=1.
\]
Its relative map is \(G=F_{Q,C/S,D}\). A power witness is an integer
\(k\in\N\) with \(S\mid Q^k\). The affine count of a finite raw orbit
segment is the number of its instructions using the complementary guard
\(D\nmid n\).
\end{definition}

The coprimality condition forces \(S\ne0\): otherwise its gcd with \(D\)
is \(|D|\ge2\). The supplied-factor conditions do not require
\(\gcd(|C|,|D|)=1\), and a power witness for \(S\) does not assert one for
\(C\).

\begin{lemma}[Exact guard and clock preservation]
\label{lem:scale}
For all \(m\in\Z\), \(r\in\N\),
\[
F(Sm)=S G(m),\qquad F^r(Sm)=S G^r(m).
\]
Consequently \(S\Z\) is invariant, and scaling is a raw-clock conjugacy
between the core state space and this sublattice.
\end{lemma}
\begin{proof}
Coprimality gives \(D\mid Sm\) if and only if \(D\mid m\), so both maps
use the same guard. If that guard holds, exact division gives
\((Sm)/D=S(m/D)\), including for signed divisors. Otherwise,
\(QSm+C=S(Qm+C/S)\), since \(S\mid C\).
This proves the one-step equality. Induction on \(r\), starting with the
identity at \(r=0\), proves the iterate equality. Multiplication by nonzero
\(S\) is a bijection from \(\Z\) to \(S\Z\), which gives the stated
conjugacy. This is not a conjugacy on the whole original state space.
\end{proof}

\section{Entry with a bounded affine count}

\begin{theorem}[Prepared entry]
\label{thm:entry}
Suppose \(S\mid Q^k n\), where \(k\in\N\) and \(n\in\Z\).
There is a raw time \(t\) with \(S\mid F^t(n)\), and at most \(k\)
affine instructions occur before that time.
\end{theorem}
\begin{proof}
Use the lexicographic order on \((k,|n|)\in\N^2\), with the exponent
first. If \(S\mid n\), choose \(t=0\). If \(k=0\), prepared divisibility
already says \(S\mid n\), so this also settles the exponent-zero case.
Assume instead that \(k>0\) and \(S\nmid n\).

If \(D\mid n\), then \(n\ne0\), since every integer divides zero.
Write \(n=Dn'\). The hypothesis says \(S\mid Q^k Dn'\);
coprimality cancels \(D\) and yields \(S\mid Q^k n'\).
Moreover \(|n'|=|n|/|D|<|n|\), because \(|D|\ge2\).
Thus the exponent stays fixed while the second rank coordinate decreases.
The current instruction is division and contributes no affine count.

If \(D\nmid n\), the next state is \(n'=Qn+C\). Both
\(Q^k n\) and \(Q^{k-1}C\) are divisible by \(S\), hence
\[
S\mid Q^{k-1}n'=Q^k n+Q^{k-1}C.
\]
The exponent decreases, irrespective of the possibly increased magnitude.
This instruction contributes one affine event; induction provides at most
\(k-1\) more. Each case has one strictly smaller child, preserves the
prepared divisibility condition, and reconstructs the original time by
prepending its one raw instruction. The well-founded induction therefore
terminates and proves the bound.
\end{proof}

\begin{theorem}[Common raw tail]
\label{thm:tail}
Assume \(S\mid Q^k\). Every signed seed \(n\) has a common entry time
\(t\) and an integer \(m\) such that
\[
F^t(n)=Sm,\qquad F^{t+r}(n)=S G^r(m)\quad(r\in\N).
\]
The prefix contains at most \(k\) affine instructions.
\end{theorem}
\begin{proof}
The power witness implies \(S\mid Q^k n\). Theorem~\ref{thm:entry}
gives \(t\) and an exact factorization \(F^t(n)=Sm\), with the required
count. For every \(r\), the iteration identity and Lemma~\ref{lem:scale}
give
\[
F^{t+r}(n)=F^r(F^t(n))=F^r(Sm)=S G^r(m).
\]
The same \(t,m\) work for all future raw times. They can depend on the
seed; no uniform all-seed entry time is asserted. The affine count is not
a raw-time bound.
\end{proof}

\begin{lemma}[A power from prime support]
\label{lem:support}
If every prime divisor of \(|S|\) divides \(|Q|\), then
\(S\mid Q^{|S|}\). This witness is sufficient and nonoptimal.
\end{lemma}
\begin{proof}
It suffices to prove that positive natural numbers \(a,q\) with supported
primes satisfy \(a\mid q^a\); the same proof allows \(q=0\).
For \(a=1\), divisibility is immediate. For \(a>1\), choose a prime
\(p\mid a\) and put \(b=a/p\). Such a prime exists by taking the least
divisor greater than one: if that divisor were composite, a smaller
nontrivial divisor would also divide \(a\).
We have \(a=pb\), \(b\ge1\), and \(b<a\). Every prime divisor of \(b\)
also divides \(a\), so induction gives \(b\mid q^b\). Since \(p\mid q\),
multiplication gives \(a\mid q^{b+1}\). The inequalities
\(p\ge2\) and \(b\ge1\) yield \(b+1\le pb=a\), and divisibility
persists upon increasing the exponent. Thus \(a\mid q^a\).
Apply this to \(a=|S|\), \(q=|Q|\), then use absolute-value divisibility
to recover the signed integer statement. The exponent-zero unit case also
avoids any exception involving \(Q=0\).
\end{proof}

\section{Arbitrary targets and positive domains}

\begin{theorem}[Universal target recovery]
\label{thm:targets}
Let \(S\mid Q^k\), and let \(P\) be any predicate on \(\Z\). Then
\[
\bigl(\forall m\in\Z\ \exists r\in\N:\ P(SG^r(m))\bigr)
\quad\Longleftrightarrow\quad
\bigl(\forall n\in\Z\ \exists t\in\N:\ P(F^t(n))\bigr).
\]
There is no invariance condition on \(P\).
If \(S>0\), \(Q\ge0\), \(C>0\), and \(D>0\), the same equivalence
holds with both seed quantifiers restricted to positive integers.
\end{theorem}
\begin{proof}
For the forward implication, choose the common tail \(t,m\) for the given
original seed. A core witness \(r\) recovers the target at original time
\(t+r\). For the reverse implication, apply original capture to the
original seed \(Sm\); Lemma~\ref{lem:scale} supplies the identical raw
clock in the core. No target value is changed or transported through an
uncontrolled intermediate state.

Under the positive hypotheses, every original raw iterate of a positive
seed is positive: exact positive division has a positive integer quotient,
and \(Qn+C>0\) at an affine instruction. At entry, \(Sm>0\) and \(S>0\)
imply \(m>0\). Conversely \(m>0\) gives a positive original seed \(Sm\).
The preceding two implications therefore apply with the positive quantifiers.
Also \(C/S>0\): exact divisibility reconstructs \(C=S(C/S)\), and a
nonpositive quotient would contradict \(C>0\). Thus the normalized
parameters stay in the positive class.
\end{proof}

The theorem is an all-seed equivalence for the pulled-back core target.
It is not a pointwise equivalence between the same original and normalized
seed, and it does not prove capture for a core map whose capture problem
remains unresolved.

\section{Every phase and the primitive raw period}

\begin{theorem}[Complete cycle recovery]
\label{thm:cycles}
Suppose \(S\mid Q^k\). Every periodic original state belongs to \(S\Z\).
For every \(m\in\Z\) and every \(r\in\N\),
\[
F^r(Sm)=Sm\quad\Longleftrightarrow\quad G^r(m)=m.
\]
Every original cycle is therefore the scaled cycle of one core state, and
all of its raw phases and its primitive raw period agree under scaling.
Every core cycle lifts in the reverse direction.
\end{theorem}
\begin{proof}
Let \(F^h(n)=n\) with \(h>0\), and choose the common entry \(t,m\) from
Theorem~\ref{thm:tail}. The multiple \(H=h(t+1)\) satisfies \(H\ge t\),
and repetition of the return gives \(F^H(n)=n\). The common-tail identity
at \(H-t\) then gives
\[
n=F^H(n)=S G^{H-t}(m).
\]
Thus the periodic state itself, not just one later phase, belongs to the
lattice. Write it as \(Sn_0\). Lemma~\ref{lem:scale} now reconstructs every
phase from the one state \(n_0\), beginning at raw time zero.
For each \(r\), injectivity of multiplication by \(S\ne0\) makes
\(S G^r(n_0)=Sn_0\) equivalent to \(G^r(n_0)=n_0\).
The sets of positive return times are identical, so their least elements
are identical. This proves primitive-period preservation against every
shorter positive raw return. The same iterate equality lifts a core cycle
and all its phases. Zero, which has raw period one, is included.
\end{proof}

\begin{lemma}[Exact starting cutoff]
\label{lem:ceiling}
For natural numbers \(S,H,n\) with \(S>0\),
\[
H\le Sn\quad\Longleftrightarrow\quad
\left\lfloor\frac{H+S-1}{S}\right\rfloor\le n.
\]
\end{lemma}
\begin{proof}
Put \(a=H+S-1\). Integer division gives
\(\lfloor a/S\rfloor\le n\) if and only if \(a<S(n+1)\).
Since \(H+S\ge1\), this strict inequality is exactly \(H\le Sn\).
The argument includes \(H=0\). Applied to a positive cycle and a positive
scale, its minimum scales by \(S\), so the original starting cutoff becomes
the displayed core cutoff. No cycle enumeration or exhaustion is inferred
from this conversion alone.
\end{proof}

\section{Gcd-only construction and its exact recovery}

\begin{theorem}[Constructed maximal supported factor]
\label{thm:peel}
For \(C\ne0\), repeatedly remove \(g=\gcd(h,|Q|)\) from
\(h\), starting with \(h=|C|\), and multiply the removed factors into
\(\Sigma\), starting with \(\Sigma=1\). Stop at \(g=1\). The result obeys
\[
\Sigma,h>0,\qquad \Sigma h=|C|,\qquad \gcd(h,|Q|)=1,
\]
and has a power witness \(\Sigma\mid Q^k\) for some \(k\le|C|\).
Both \(\Sigma\) and \(h\) are at most \(|C|\).
For every positive divisor \(a\mid|C|\), its primes all divide \(|Q|\)
if and only if \(a\mid\Sigma\).
If \(\gcd(|C|,|D|)=1\), the computed \(\Sigma\) supplies the hypotheses
of Theorems~\ref{thm:tail}, \ref{thm:targets}, and \ref{thm:cycles}.
The relative offset has magnitude \(h\), hence is coprime in magnitude to
\(Q\). If \(C>0\) and \(\Sigma>1\), it is strictly smaller than \(C\).
\end{theorem}
\begin{proof}
A nonterminal gcd divides positive \(h\), so it is positive and, being
unequal to one, is at least two. Thus \(h/g<h\), proving termination in
the positive natural class. At a terminal node, the certificate is
\((1,h)\), with exponent zero. At a nonterminal node, let the child's
certificate be \((\Sigma',h')\) with exponent \(k'\le h/g\).
The new pair \((g\Sigma',h')\) reconstructs \(h\). Because \(g\mid|Q|\)
and \(\Sigma'\mid |Q|^{k'}\), it has exponent \(k'+1\).
The inequalities \(g\ge2\) and \(h/g<h\) imply
\(k'+1\le h/g+1\le h\). The terminal core is coprime to \(|Q|\).
Induction proves the product, positivity and exponent bounds. Positivity
of both factors in the product bounds each by the original \(|C|\).

A prime dividing \(\Sigma\) divides \(|Q|\), since it divides a power
of \(|Q|\); this follows by induction on the power using primality.
Conversely, let \(a\mid|C|=\Sigma h\) have supported primes.
Lemma~\ref{lem:support} gives \(a\mid|Q|^a\), while
\(\gcd(h,|Q|)=1\) implies \(\gcd(h,|Q|^a)=1\).
Thus \(\gcd(a,h)=1\), and cancellation from \(a\mid\Sigma h\) gives
\(a\mid\Sigma\). This proves maximality in the divisibility order.

The factor \(\Sigma\) divides \(|C|\), hence divides \(C\).
If \(C\) is coprime to \(D\), so is this divisor. The constructed power
witness therefore applies to the same original map, not to a separately
chosen scale. Exact division gives
\(\Sigma(C/\Sigma)=C\); taking absolute values and cancelling positive
\(\Sigma\) against \(\Sigma h=|C|\) gives \(|C/\Sigma|=h\).
This proves core coprimality and compatibility with the relative parameters.
For a positive original offset and \(\Sigma>1\), positive exact division
strictly decreases the offset. If \(\Sigma=1\), the core was already
coprime to the multiplier, and there is no strict decrease.
The case \(Q=0\) obeys the same certificates: for \(|C|>1\) the first
gcd removes the whole offset, leaving core one and exponent one.
The constructor is also defined at zero offset by \((1,0)\), but the
positive certificate and this theorem exclude that fallback.
\end{proof}

\section{Examples and exact arithmetic}

For \((Q,C,D)=(11,55,2)\), gcd removal gives \(\Sigma=11\), core offset
\(5\), and exponent one. The seed \(1\) enters as
\(1\mapsto66=11\cdot6\). Requiring the whole offset to divide a power
would fail: every power of \(11\) is congruent to one modulo \(5\), while
\(5\mid55\). For \((3,6,2)\), the supplied factor \(3\) gives core offset
\(2\) although the full offset is not coprime to the divisor. The generic
supplied-factor theorem applies; the stronger full-offset criterion is only
sufficient for automatic admissibility.

For \((181,181^{20},2)\), extraction gives core \((181,1,2)\).
The exact arithmetic
\[
181\cdot27+1=611\cdot2^3,\qquad
181\cdot611+1=27\cdot2^{12}
\]
and
\[
181\cdot35+1=99\cdot2^6,\qquad
181\cdot99+1=35\cdot2^9
\]
produces the two odd-state pairs.
Their respective raw words are
\[
O D^3 O D^{12},\qquad O D^6 O D^9.
\]
Here \(O\) means affine and \(D\) labels a division instruction. Each word
has raw length \(17\). Exact divisibility and odd endpoints verify every guard.
Within a division chain the states are distinct, and the two chains have
different odd parts; consequently all seventeen raw cycle states are
distinct. Both periods are primitive, and the minima are respectively
\(27\) and \(35\). Scaling by the odd factor \(181^{20}\) preserves
every raw edge and multiplies these minima by that factor.

The exact Go replay uses signed arbitrary-precision integers. It checks the
constructed power certificate, a common entry, the affine budget, and
sixteen normalized tail edges, recomputing both complementary guards.
The fixed box has \(-7\le Q\le7\), \(1\le|C|\le32\),
\(2\le|D|\le8\), full-offset coprimality, and \(-64\le n\le64\).
For divisor magnitudes \(2,3,4,5,6,7,8\), the counts of coprime positive
offsets through \(32\) are \(16,22,16,26,11,28,16\), summing to \(135\).
The box therefore has exactly
\[
15\cdot4\cdot135\cdot129=1{,}044{,}900
\]
seed/parameter certificates, not distinct orbits. Additional controls cover
supplied factors, signed divisions, zero multipliers, both scaled raw cycles,
large signed seeds, ceiling boundaries and invalid inputs.

The replay accepts nonzero offsets and positive supplied factors;
the signed-factor and zero-offset entry theorems have wider domains.
It restricts coefficients and supplied raw or core seeds to 512 bits,
\(|Q|\le181\), \(2\le|D|\le32\), and entry prefixes to 4096 raw steps.
These limits do not restrict the formal statements. A resource refusal is
not nonabsorption, divergence or search exhaustion. The replay is not a
complete core-cycle search. Zero allocations are measured only in four
warmed configurations with seed \(27\).

\section{Formal statements and commands}

The module correspondence is as follows.
\begin{itemize}
\item \path{lean/SupportedOffsetAbsorption.lean} contains the supplied-factor
entry rank, common tail, target equivalences, return-time reflection,
primitive-cycle recovery and exact cutoff.
\item \path{lean/SupportedOffsetFactor.lean} proves gcd extraction.
Its certificates retain products, support, power bounds and maximality.
\item \path{lean/SupportedOffsetExtracted.lean} uses the actual constructed
factor in the raw recovery and proves the core's multiplier coprimality.
\item The proper-factor obstruction and raw-period controls are in
\path{lean/SupportedOffsetControls.lean}. They exclude every shorter return
of both core cycles and lift their periods to the large offset.
\end{itemize}
The core parameter class, zero-offset exception, signed guards and positive
capture hypotheses are retained at the corresponding use sites.
From the repository root, use \texttt{make lean}, the \texttt{go} profile
with \path{QCD.Scaling.002/go/SupportedOffsetAbsorption.go}, or the
\texttt{tex} profile with this document as \texttt{TEXFILE}.

\bibliographystyle{amsalpha}
\bibliography{SupportedOffsetAbsorption}

\end{document}
